\section{Erfahrungen und Erkenntnisse aus dem Modul Vernetzte Systeme}\label{Semesterrückblick}

\subsection{Erlernte Kompetenzen und Kenntnisse im Semester }
Das Modul „Vernetzte Systeme“ hat uns in diesem Semester eine Vielzahl von Konzepten, Technologien und praktischen Anwendungen nähergebracht. Der Fokus lag darauf, Systeme zu entwerfen, die über mehrere Rechner oder Prozesse verteilt sind und effizient miteinander kommunizieren können und diese anschließend zu testen mit Werkzeugen erhöhte Last simulieren. Dazu haben wir verschiedene Werkzeuge und Technologien kennengelernt, darunter Docker, TCP/IP, das ISO/OSI-Modell, SQL-Datenbanken, HTTP-Protokolle und Webprogrammierung, sowie diverse Lasttest-Tools.
\subsubsection{Vernetzte Systeme und das Schichtenmodell }
Ein zentraler Bestandteil des Moduls war das Verständnis der Grundlagen vernetzter Systeme. Wir haben gelernt, dass moderne Softwarelösungen nicht mehr auf einzelnen Servern laufen, sondern oft auf verteilten Systemen basieren. Diese Systeme bestehen aus mehreren Komponenten, die miteinander kommunizieren und Daten austauschen müssen.
\subsubsection{Webprogrammierung und HTTP }
Da moderne Anwendungen oft über das Web kommunizieren, haben wir das HTTP-Protokoll detailliert behandelt. Wir haben uns mit HTTP, Cookies und Sessions beschäftigt und gelernt, wie Webserver und Clients miteinander kommunizieren.

Mit HTML, CSS und JavaScript konnten wir eigene Webanwendungen erstellen, die über AJAX dynamisch Daten nachladen. Dies war besonders wichtig für Anwendungen, die in Echtzeit mit dem Server kommunizieren müssen – ein Bereich, in dem WebSockets eine große Rolle spielen.

Dafür haben wir uns das ISO/OSI-Schichtenmodell sowie den TCP/IP-Stack angesehen. Diese Modelle helfen dabei, Netzwerke strukturiert zu verstehen und Protokolle gezielt einzusetzen. Besonders wichtig war die Transportschicht (TCP/UDP), da sie für die zuverlässige Kommunikation zwischen Servern und Clients sorgt.
\subsubsection{Linux-Umgebung und Posix-Standards }
Da viele vernetzte Systeme auf Linux-Servern betrieben werden, haben wir uns intensiv mit Linux-Befehlen und der Arbeit in der Bash-Shell beschäftigt. Wir haben uns die Umgebung bereits im vorherigen Semester mehrmals aufgesetzt und gelernt mit den Tools umzugehen, um Daten effizient zu verarbeiten. Somit konnten wir alle mit der gleichen Umgebung arbeiten, unabhängig vom Betriebssystem des privaten Endgeräts. Ein weiteres wichtiges Konzept war Posix (Portable Operating System Interface). Posix stellt sicher, dass Software auf verschiedenen Unix-basierten Systemen lauffähig bleibt.
Da Docker-Container ebenfalls auf Posix basieren, erleichtert dies die Portabilität von Anwendungen.
\subsubsection{Datenbanken und Datenverwaltung}
Ein weiterer großer Themenbereich war der Umgang mit SQL-Datenbanken (MariaDB). Wir haben bereits im vorherigen Semester gelernt, wie man Datenbanken anlegt, Tabellen erstellt, Daten einfügt, aktualisiert und abfragt. Auf dieses Wissen haben wir dieses Semester aufgebaut und unsere Kenntnisse verbessert.

Zusätzlich haben wir untersucht, wie sich Datenbank-Architekturen in vernetzten Systemen optimieren lassen. Hierzu gehörten Konzepte wie Datenbank-Sharding und Caching mit Redis.
\subsubsection{Docker und Containerisierung }
Ein zentrales Thema des Semesters war die Container-Technologie Docker. Wir haben gelernt, inwiefern Docker uns hilft, Anwendungen und ihre Abhängigkeiten isoliert in Containern bereitzustellen. Dies erleichtert die Skalierung und den Betrieb verteilter Systeme erheblich. In unserem Projekt haben wir Docker intensiv genutzt, um unsere Webanwendung in einer vernetzten Umgebung aufzubauen. Die Grundlage die wir genutzt haben wurde freundlicherweise von unserem Dozenten zur Verfügung gestellt. In diesem Verzeichnis sind die Fundamente, auf denen wir unser eigenes Projekt aufbauen konnten, enthalten. Durch das bereits funktionsfähige Vernetzte System, wie z.B build-all.sh oder die bin/ und context/ Verzeichnisse der einzelnen Docker-Container konnten wir uns nach herzenslust ausprobieren und immer wieder von Neuem beginnen, wenn wir uns in eine Sackgasse gearbeitet haben.

\subsubsection{Lasttest}

